% Part: incompleteness
% Chapter: introduction
% Section: definitions

\documentclass[../../../include/open-logic-section]{subfiles}

\begin{document}

\olfileid[hi]{inc}{int}{def}

\olsection{परिभाषाएँ}

गणित को औपचारिक बनाने और यह दिखाने की हिल्बर्ट की परियोजना पूरी करने के लिए
कि ऐसा औपचारिकीकरण संगत और पूर्ण है, पहला कार्य किसी भाषा, तार्किक ढाँचे और
अभिगृहीत-तंत्र का चयन करना होगा। अपने प्रयोजनों के लिए मान लें कि गणित को
किसी प्रथम-कोटि भाषा में औपचारिक बनाया जा सकता है; अर्थात् !!{constant},
!!{function} और !!{predicate} के कुछ ऐसे समुच्चय हैं जो प्रथम-कोटि
तर्कशास्त्र के संयोजकों और परिमाणकों के साथ मिलकर हमें गणित के कथन व्यक्त
करने देते हैं। अधिकांश लोग सहमत हैं कि ऐसी भाषा विद्यमान है: समुच्चय सिद्धांत
की भाषा, जिसमें $\in$ एकमात्र अतार्किक चिह्न है। पहली दृष्टि में यह दावा बहुत
अविश्वसनीय लगता है कि इतनी सरल भाषा इतनी अभिव्यंजक है; और यह स्थापित करने में
बहुत कार्य लगा कि लगभग संपूर्ण गणित को इस अत्यंत मित शब्दावली में व्यक्त किया
जा सकता है। बात सरल रखने के लिए अभी अपनी चर्चा को अंकगणित, अर्थात् गणित के
उस भाग तक सीमित करें जो केवल प्राकृतिक संख्याओं~$\Nat$ से संबंधित है।
अंकगणित के तथ्यों को व्यक्त करने की स्वाभाविक भाषा~$\Lang L_A$ है। $\Lang
L_A$ में एक द्विस्थानी !!{predicate}~$<$, एक !!{constant}~$\Obj 0$, एक
एकस्थानी !!{function}~$\prime$, और दो द्विस्थानी !!{function}~$+$ तथा
~$\times$ हैं।

\begin{defn}
!!{sentence} का समुच्चय~$\Gamma$ एक \emph{सिद्धांत} है यदि वह तार्किक
फलन के अंतर्गत संवृत है; अर्थात् यदि $\Gamma = \Setabs{!A}{\Gamma \Entails
!A}$।
\end{defn}

सिद्धांतों को निर्दिष्ट करने की दो सरल विधियाँ हैं। पहली है किसी
!!{structure} में सत्य !!{sentence} के समुच्चय के रूप में। उदाहरण के लिए,
$\Lang L_A$ के उस !!{structure} पर विचार करें जिसमें !!{domain}~$\Nat$ है और
सभी अतार्किक चिह्नों की व्याख्या अपेक्षित ढंग से की गई है।

\begin{defn}\ollabel{def:standard-model}
\emph{अंकगणित का मानक प्रतिरूप} निम्न प्रकार परिभाषित !!{structure}~$\Struct
N$ है:
\begin{enumerate}
\item $\Domain N = \Nat$
\item $\Assign{\Obj 0}{N} = 0$
\item $\Assign{\Obj \prime}{N}(n) = n + 1$, प्रत्येक $n \in \Nat$ के लिए
\item $\Assign{\Obj +}{N}(n, m) = n + m$, प्रत्येक $n, m \in \Nat$ के लिए
\item $\Assign{\Obj \times}{N}(n, m) = n\cdot m$, प्रत्येक $n, m \in \Nat$ के लिए
\item $\Assign{\Obj <}{N} = \Setabs{\tuple{n, m}}{n \in \Nat, m \in
  \Nat, n < m}$
\end{enumerate}
\end{defn}

`$\times$' और `$\cdot$' के बीच अंतर पर ध्यान दें: $\times$ अंकगणित की भाषा
का एक चिह्न है। निस्संदेह, हमने इसे इसलिए चुना है कि यह हमें गुणन की याद
दिलाए; किंतु $\times$ गुणन संक्रिया नहीं, बल्कि एक द्विस्थानी फलन-चिह्न
(औपचारिक रूप से $\Obj f^2_1$) है। इसके विपरीत, $\cdot$ साधारण गुणन
फलन~\emph{है}। $n \cdot m$ जैसी अभिव्यक्ति में हमारा आशय संख्याओं $n$ और
~$m$ के गुणनफल से है; $x \times y$ जैसी अभिव्यक्ति में हम अंकगणित की भाषा
के एक पद की बात कर रहे होते हैं। मानक प्रतिरूप में फलन-चिह्न~$\times$ की
व्याख्या प्राकृतिक संख्याओं पर फलन $\cdot$ के रूप में होती है। योग के लिए हम
$+$ का उपयोग अंकगणित की भाषा के फलन-चिह्न और प्राकृतिक संख्याओं पर योग-फलन,
दोनों के रूप में करते हैं। यहाँ आशय संदर्भ से निर्धारित करना होता है।

\begin{defn}
\emph{सत्य अंकगणित} का सिद्धांत अंकगणित के मानक प्रतिरूप में संतुष्ट
!!{sentence} का समुच्चय है; अर्थात्
\[
\Th{TA} = \Setabs{!A}{\Sat{N}{!A}}.
\]
\end{defn}

$\Th{TA}$ एक सिद्धांत है, क्योंकि जब भी $\Th{TA} \Entails !A$ होता है, तब
$!A$ प्रत्येक ऐसे !!{structure} में संतुष्ट होता है जो~$\Th{TA}$ को संतुष्ट
करता है। चूँकि $\Sat{N}{\Th{TA}}$, अतः~$\Sat{N}{!A}$ और इसलिए $!A \in
\Th{TA}$।

सिद्धांत~$\Gamma$ को निर्दिष्ट करने की दूसरी विधि उसे उन !!{sentence} का
समुच्चय मानना है जो वाक्यों के किसी समुच्चय~$\Gamma_0$ से तार्किक रूप से
फलित होते हैं। तब $\Gamma$, तार्किक फलन के अंतर्गत $\Gamma_0$ का ``संवरण''
है। इस प्रकार किसी सिद्धांत को निर्दिष्ट करना तभी रोचक है जब $\Gamma_0$ को
स्पष्ट रूप से दिया गया हो; उदाहरणतः~$\Gamma_0$ के !!{element} सूचीबद्ध हों।
कम-से-कम $\Gamma_0$ निर्णेय होना चाहिए; अर्थात् यह जाँचने के लिए कोई संगणनीय
परीक्षण होना चाहिए कि कोई !!{sentence},~$\Gamma_0$ का अवयव है या नहीं।
हम~$\Gamma_0$ के !!{sentence} को~$\Gamma$ के \emph{अभिगृहीत} कहते हैं, और
कहते हैं कि $\Gamma$,~$\Gamma_0$ द्वारा \emph{अभिगृहीतीकृत} है।

\begin{defn}
सिद्धांत~$\Gamma$,~$\Gamma_0$ द्वारा \emph{अभिगृहीतीकृत} है, यदि और केवल यदि
\[
\Gamma = \Setabs{!A}{\Gamma_0 \Entails !A}
\]
\end{defn}

\begin{defn}
निम्न वाक्यों द्वारा अभिगृहीतीकृत सिद्धांत $\Th{Q}$ ``रॉबिन्सन का
$\Th{Q}$'' कहलाता है और अंकगणित का एक बहुत सरल सिद्धांत है।
\begin{align*}
& \lforall[x][\lforall[y][(\eq[x'][y'] \lif \eq[x][y])]] \tag{$!Q_1$}\\
& \lforall[x][\eq/[\Obj 0][x']] \tag{$!Q_2$}\\
& \lforall[x][(\eq[x][\Obj 0] \lor \lexists[y][\eq[x][y']])] \tag{$!Q_3$}\\
& \lforall[x][\eq[(x + \Obj 0)][x]] \tag{$!Q_4$}\\
& \lforall[x][\lforall[y][\eq[(x + y')][(x + y)']]] \tag{$!Q_5$}\\
& \lforall[x][\eq[(x \times \Obj 0)][\Obj 0]] \tag{$!Q_6$}\\
& \lforall[x][\lforall[y][\eq[(x \times y')][((x \times y) + x)]]] \tag{$!Q_7$}\\
& \lforall[x][\lforall[y][(x < y \liff \lexists[z][\eq[(z' + x)][y]])]] \tag{$!Q_8$}
\end{align*}
!!{sentence} का समुच्चय $\{!Q_1, \dots, !Q_8\}$, $\Th{Q}$ के अभिगृहीत हैं;
अतः $\Th{Q}$ उन सभी !!{sentence} से बना है जो उनसे तार्किक रूप से फलित होते
हैं:
\[
\Th{Q} = \Setabs{!A}{\{!Q_1, \dots, !Q_8\} \Entails !A}.
\]
\end{defn}

\begin{defn}
मान लें कि $!A(x)$, $\Lang L_A$ में मुक्त चरों~$x$ और $y_1$, \dots,
$y_n$ वाला !!a{formula} है। तब निम्न रूप का कोई भी !!{sentence}
\[
\lforall[y_1][\dots\lforall[y_n][((!A(\Obj 0) \land \lforall[x][(!A(x)
\lif !A(x'))]) \lif \lforall[x][!A(x)])]]
\]
\emph{आगमन-योजना} का एक उदाहरण है।

\emph{पियानो अंकगणित}~$\Th{PA}$ वह सिद्धांत है जो $\Th{Q}$ के अभिगृहीतों और
आगमन-योजना के सभी उदाहरणों द्वारा अभिगृहीतीकृत है।
\end{defn}

\begin{explain}
आगमन-योजना का प्रत्येक उदाहरण~$\Struct{N}$ में सत्य है। इसे देखना तब सबसे
सरल है जब !!{formula}~$!A$ में केवल एक मुक्त !!{variable}~$x$ हो। तब
$!A(x)$,~$\Struct{N}$ में~$\Nat$ का उपसमुच्चय~$X_{!A}$ परिभाषित करता है।
$X_{!A}$ उन सभी~$n \in \Nat$ का समुच्चय है जिनके लिए $s(x) = n$ होने पर
$\Sat{N}{!A(x)}[s]$। आगमन-योजना का संबंधित उदाहरण है
\[
((!A(\Obj 0) \land \lforall[x][(!A(x) \lif !A(x'))]) \lif 
  \lforall[x][!A(x)]).
\]
यदि इसका पूर्वपक्ष~$\Struct{N}$ में सत्य है, तो $0 \in X_{!A}$ और जब भी
$n \in X_{!A}$, तब $n+1$ भी उसमें है। चूँकि $0 \in X_{!A}$, हमें $1 \in
X_{!A}$ मिलता है। $1 \in X_{!A}$ से $2 \in X_{!A}$ मिलता है, और यही क्रम
चलता रहता है। अतः प्रत्येक $n \in \Nat$ के लिए $n \in X_{!A}$। इसका अर्थ
है कि $\lforall[x][!A(x)]$,~$\Struct{N}$ में संतुष्ट है।
\end{explain}

दोनों $\Th{Q}$ और $\Th{PA}$ अभिगृहीतीकृत सिद्धांत हैं। बड़ा प्रश्न है कि वे
कितने प्रबल हैं। उदाहरणतः, क्या $\Th{PA}$,~$\Nat$ के विषय में वे सभी सत्य
सिद्ध कर सकता है जिन्हें~$\Lang L_A$ में व्यक्त किया जा सकता है? विशेषतः,
क्या $\Th{PA}$ के अभिगृहीत उन सभी प्रश्नों का निपटारा करते हैं जिन्हें~$\Lang
L_A$ में सूत्रित किया जा सकता है?

इसे पूछने का दूसरा ढंग है: क्या $\Th{PA} = \Th{TA}$? स्पष्टतः $\Th{TA}$,
~$\Nat$ के विषय में सभी सत्यों को सिद्ध करता है (अर्थात् उन्हें समाहित करता
है), और वह~$\Lang L_A$ में सूत्रित किए जा सकने वाले सभी प्रश्नों का निपटारा
करता है; क्योंकि यदि $!A$, $\Lang L_A$ में !!a{sentence} है, तो या तो
$\Sat{N}{!A}$ अथवा $\Sat{N}{\lnot !A}$, और इसलिए या तो $\Th{TA} \Entails
!A$ अथवा $\Th{TA} \Entails \lnot !A$। ऐसे सिद्धांत को \emph{!!{complete}}
कहें।

\begin{defn}
सिद्धांत $\Gamma$, \emph{!!{complete}} है यदि और केवल यदि उसकी भाषा के
प्रत्येक !!{sentence}~$!A$ के लिए या तो $\Gamma \Entails !A$ अथवा $\Gamma
\Entails \lnot !A$।
\end{defn}

\begin{explain}
पूर्णता प्रमेय के अनुसार, $\Gamma \Entails !A$ यदि और केवल यदि $\Gamma
\Proves !A$; अतः $\Gamma$ पूर्ण है यदि और केवल यदि उसकी भाषा के प्रत्येक
!!{sentence}~$!A$ के लिए या तो $\Gamma \Proves !A$ अथवा $\Gamma \Proves
\lnot !A$।
\end{explain}

इससे एक और प्रश्न उठता है: क्या कोई ऐसी संगणकीय प्रक्रिया है जिससे हम जाँच
सकें कि कोई !!a{sentence},~$\Th{TA}$ में, $\Th{PA}$ में, अथवा केवल~$\Th{Q}$
में ही है या नहीं? कोई समुच्चय (उदाहरणतः !!{sentence} का समुच्चय) कब
!!{decidable} है, इसे परिभाषित करके हम प्रश्न को अधिक सटीक बना सकते हैं।

\begin{defn}
समुच्चय $X$, \emph{!!{decidable}} है यदि और केवल यदि कोई ऐसी संगणकीय
प्रक्रिया है जो आगत~$x$ पर, $x \in X$ होने पर~$1$ और अन्यथा $0$ लौटाती है।
\end{defn}

अतः हमारा प्रश्न बनता है: क्या $\Th{TA}$ ($\Th{PA}$, $\Th{Q}$) !!{decidable}
है?

इन सभी प्रश्नों का उत्तर होगा: नहीं। इनमें से कोई सिद्धांत निर्णेय नहीं है।
किंतु यह परिघटना केवल इन्हीं विशिष्ट सिद्धांतों तक सीमित नहीं। वास्तव में
कुछ शर्तें पूरी करने वाला \emph{कोई भी} सिद्धांत इन्हीं परिणामों के अधीन है।
इनमें एक शर्त, जिसे $\Th{Q}$ और $\Th{PA}$ पूरा करते हैं, यह है कि वे
अभिगृहीतों के !!a{decidable} समुच्चय द्वारा अभिगृहीतीकृत हैं।

\begin{defn}
कोई सिद्धांत \emph{!!{axiomatizable}} है यदि वह अभिगृहीतों के
!!a{decidable} समुच्चय द्वारा अभिगृहीतीकृत है।
\end{defn}

\begin{ex}
!!{sentence} के किसी परिमित समुच्चय द्वारा अभिगृहीतीकृत प्रत्येक सिद्धांत
!!{axiomatizable} है, क्योंकि प्रत्येक परिमित समुच्चय !!{decidable} होता है।
इसलिए, उदाहरणतः, $\Th{Q}$ !!{axiomatizable} है।

~$\Th{PA}$ जैसे योजनाबद्ध रूप से अभिगृहीतीकृत सिद्धांत भी
!!{axiomatizable} हैं। यह जाँचने के लिए कि $!B$,~$\Th{PA}$ के अभिगृहीतों में
है या नहीं—अर्थात् उस फलन $\Char{X}$ की गणना करने के लिए जिसमें यदि $!B$,
~$\Th{PA}$ का अभिगृहीत है तो $\Char{X}(!B) = 1$ और अन्यथा $= 0$—हम यह कर
सकते हैं। पहले जाँचें कि $!B$,~$\Th{Q}$ के अभिगृहीतों में से एक है या नहीं।
यदि है, तो उत्तर ``हाँ'' और $\Char{X}(!B) = 1$ है। यदि नहीं, तो जाँचें कि
वह आगमन-योजना का उदाहरण है या नहीं। यह व्यवस्थित ढंग से किया जा सकता है; इस
स्थिति में शायद निम्न विधि से इसे देखना सबसे सरल है। आगमन-योजना का प्रत्येक
उदाहरण अनेक सार्वत्रिक परिमाणकों से आरंभ होता है, और उसके बाद एक उप-!!{formula}
आता है जो एक प्रतिबंधात्मक है। उस प्रतिबंधात्मक का उत्तरपक्ष
$\lforall[x][!A(x, y_1, \dots, y_n)]$ है, जहाँ $x$ और $y_1$, \dots, $y_n$,
~$!A$ के सभी मुक्त चर हैं, और~$!B$ के आरंभिक परिमाणक चरों~$y_1$, \dots,
~$y_n$ को आबद्ध करते हैं। इस~$!A$ को निकालने और यह जाँचने के बाद कि उसके
मुक्त चर आरंभ के सार्वत्रिक परिमाणकों तथा~$\lforall[x]$ से आबद्ध चरों से मेल
खाते हैं, हम जाँचते हैं कि प्रतिबंधात्मक का पूर्वपक्ष निम्न से मेल खाता है:
\[
!A(\Obj 0, y_1, \dots, y_n) \land \lforall[x][(!A(x, y_1, \dots, y_n)
\lif !A(x', y_1, \dots, y_n))]
\]
यदि मेल खाता है, तो $!B$ आगमन-योजना का उदाहरण है; और यदि नहीं, तो $!B$ नहीं है।
\end{ex}

इस प्रश्न का—और कौन-से सिद्धांत पूर्ण या निर्णेय हैं, इस अधिक सामान्य प्रश्न
का—उत्तर देते समय निम्न परिभाषा पर भी विचार करना उपयोगी होगा। स्मरण करें कि
समुच्चय $X$, !!{enumerable} है यदि और केवल यदि वह रिक्त है अथवा कोई
!!a{surjective} फलन~$f \colon \Nat \to X$ है। ऐसे फलन को~$X$ का परिगणन
कहते हैं।

\begin{defn}
समुच्चय $X$ को \emph{!!{computably enumerable}} (संक्षेप में !!{c.e.}) कहते
हैं यदि और केवल यदि वह रिक्त है अथवा उसका कोई संगणनीय परिगणन है।
\end{defn}

!!{axiomatizability} के अतिरिक्त, जिन सिद्धांतों पर अपूर्णता प्रमेय लागू होते
हैं उनके लिए एक और शर्त यह होगी कि वे संगणनीय फलनों और !!{decidable} संबंधों
के मूल तथ्यों को सिद्ध करने के लिए पर्याप्त प्रबल हों। ``मूल तथ्यों'' से
हमारा आशय उन !!{sentence} से है जो प्रत्येक तर्क के लिए संगणनीय फलनों के मान
व्यक्त करते हैं। ``पर्याप्त प्रबल'' से आशय है कि संबंधित सिद्धांत इन वाक्यों
को अपने प्रमेयों में गिनते हैं। उदाहरण के लिए, एक आदर्श संगणनीय फलन—योग—पर
विचार करें। तर्कों $2$ और $3$ के लिए $+$ का मान~$5$ है, अर्थात् $2+3 = 5$।
अंकगणित की भाषा में यह व्यक्त करने वाला वाक्य कि तर्कों $2$ और $3$ के लिए
$+$ का मान~$5$ है, यह है: $(\num{2} + \num{3}) = \num{5}$। उदाहरणतः,
$\Th{Q}$ इस वाक्य को सिद्ध करता है। अधिक सामान्यतः, हम चाहेंगे कि प्रत्येक
संगणनीय फलन $f(x_1, x_2)$ के लिए~$\Lang{L_A}$ में कोई !!a{formula}
~$!A_f(x_1, x_2, y)$ ऐसा हो कि जब भी $f(n_1, n_2) = m$, तब $\Th{Q} \Proves
!A_f(\num{n_1}, \num{n_2}, \num{m})$। इस प्रकार $\Th{Q}$ सिद्ध करता है कि
तर्कों $n_1$, $n_2$ के लिए~$f$ का मान~$m$ है। वास्तव में हम उससे थोड़ा अधिक
सिद्ध करने की अपेक्षा करते हैं: तर्कों $n_1$,~$n_2$ के लिए कोई अन्य संख्या
~$f$ का मान नहीं है। !!{decidable} संबंधों के लिए भी यही बात लागू होती है।
इसे अगली दो परिभाषाएँ सटीक बनाती हैं।

\begin{defn}
!!{formula}~$!A(x_1, \dots, x_k, y)$,~$\Gamma$ में फलन $f\colon \Nat^k \to
\Nat$ को \emph{!!{represents}} करता है यदि और केवल यदि जब भी $f(n_1,
\dots, n_k) = m$, तब
\begin{enumerate}
\item $\Gamma \Proves !A(\num{n_1}, \dots, \num{n_k}, \num{m})$, और
\item $\Gamma \Proves \lforall[y](!A(\num{n_1}, \dots, \num{n_k},
y) \lif y = \num{m})$.
\end{enumerate}
\end{defn}

\begin{defn}
!!{formula}~$!A(x_1, \dots, x_k)$ संबंध $R \subseteq \Nat^k$ को
\emph{!!{represents}} करता है यदि और केवल यदि
\begin{enumerate}
\item जब भी $R(n_1, \dots, n_k)$ हो, $\Gamma \Proves !A(\num{n_1},
\dots, \num{n_k})$, और
\item जब $R(n_1, \dots, n_k)$ नहीं होता, तब $\Gamma \Proves \lnot
!A(\num{n_1}, \dots, \num{n_k})$।
\end{enumerate}
\end{defn}

अपूर्णता प्रमेय लागू होने के लिए कोई सिद्धांत तब ``पर्याप्त प्रबल'' है जब वह
सभी संगणनीय फलनों और सभी !!{decidable} संबंधों को !!{represents} करता है।
$\Th{Q}$ और उसके विस्तार यह शर्त पूरी करते हैं, किंतु इसे स्थापित करने में
हमें कुछ समय लगेगा। यह इस विषय का अतुच्छ तथ्य है कि $\Th{Q}$ किस प्रकार की
बातें सिद्ध कर सकता है; और इसे दिखाना कठिन है, क्योंकि $\Th{Q}$ में केवल कुछ
ही अभिगृहीत हैं जिनसे हमें ये सभी तथ्य सिद्ध करने होंगे। फिर भी $\Th{Q}$ एक
बहुत दुर्बल सिद्धांत है। अतः यद्यपि यह सिद्ध करना कठिन है कि $\Th{Q}$ सभी
संगणनीय फलनों को निरूपित करता है, अधिकांश रोचक सिद्धांत~$\Th{Q}$ से अधिक
प्रबल हैं, अर्थात् वे $\Th{Q}$ से अधिक सिद्ध करते हैं। और यदि $\Th{Q}$ कोई
बात सिद्ध करता है, तो प्रत्येक अधिक प्रबल सिद्धांत भी उसे सिद्ध करता है;
चूँकि $\Th{Q}$ सभी संगणनीय फलनों को निरूपित करता है, अतः प्रत्येक अधिक प्रबल
सिद्धांत भी ऐसा करता है। इसका अर्थ है कि बहुत-से रोचक सिद्धांत अपूर्णता
प्रमेयों की यह शर्त पूरी करते हैं। इसलिए हमारा कठिन कार्य सार्थक होगा, क्योंकि
वह दिखाता है कि अपूर्णता प्रमेय सिद्धांतों के एक व्यापक वर्ग पर लागू होते हैं।
निश्चित ही ``समस्त गणित'' को औपचारिक बनाने का लक्ष्य रखने वाले किसी भी
सिद्धांत को वह सब सिद्ध करना चाहिए जो $\Th{Q}$ सिद्ध करता है, क्योंकि
कम-से-कम उसे प्रारंभिक परिकलनों के परिणाम ग्रहण करने योग्य होना चाहिए। अतः
``समस्त गणित'' के सिद्धांत का प्रत्येक प्रत्याशी ऐसा सिद्धांत होगा जिस पर
अपूर्णता प्रमेय लागू होते हैं।


\end{document}
